\section{Implementation}
\subsection{PID Control on MCU}

The continuous-time PID controller designed in Simulink was discretized for implementation on the ESP32 microcontroller.
The software was developed in C++ and relies on a fixed sample time ($\Delta t$) to ensure the integral and derivative calculations remain mathematically accurate.

\subsubsection{Discrete PID Algorithm and Hardware Protection}
Because the physical plant is constrained by a 0--3.3V operating range mapped to a 0--100\% PWM duty cycle, several engineering protections were implemented in the firmware:

\begin{itemize}
    \item \textbf{PID output limiting:} The \texttt{pid\_out} variable is constrained to the physical limit of 3.3V.
    \item \textbf{Anti-derivative kick:} The first derivative calculation will spike to infinity causing unpredictable behavior, this is handled by implementing a \texttt{first\_run} flag variable that sets the previous error equal to the current error when the derivative is calculated for the first time.
    \item \textbf{Anti-integral windup:} The integral term does not get updated if the system demands an output higher than the ESP's limit and the error is not negative. The inverse case is handled in the same way.
\end{itemize}




